\section{Experimental setup}
\label{sec:experimental-setup}

\paragraph*{Protocol.}
Every \texttt{.pog} file is translated twice, by \beer{} (revision \texttt{83fda63}, built with Lean~4.27.0) and by the ppTrans translator of Atelier~B~24.04.2.
Both emit exactly one \smtinl{check-sat} per goal, in the order the goals appear in the file, and each goal becomes a self-contained script solved in its own process---no incremental mode, no state shared between goals---so answer \(k\) on one side and answer \(k\) on the other describe the same obligation.
Files where the two goal counts differ are excluded from the join rather than silently misaligned.
Each script is discharged by cvc5~1.3.4 with model-based quantifier instantiation (\texttt{-{}-mbqi}; \Cref{sec:ar-smt}) and a per-\smtinl{check-sat} budget of \(5\)\,s (\texttt{-{}-tlimit-per=5000}), the configuration under which \beer{}'s scripts are intended to run (\Cref{sec:beer-po-running-example}).
The budget is soft---cvc5 checks it between phases---so each solver process is additionally killed at a \(300\)\,s wall; the encoders themselves run under a \(1\,800\)\,s budget per file.
\beer{}'s shipped prelude requests unsat cores, which ppTrans never does and which costs cvc5 real time, so the option is stripped and both encodings are solved under identical solver settings.

\paragraph*{Hardware.}
The campaign ran on the \texttt{gros} cluster of Grid'5000\footnote{\url{https://www.grid5000.fr}} (Nancy): identical nodes with one Intel Xeon Gold 5220 processor, running one solver instance per physical core.
All timings quoted below are from this configuration.

\paragraph*{Scored population.}
Three files had their sweep stopped partway and contribute only the goals actually measured; in total, \bmFiles{} of the corpus's 5\,421 files contribute at least one scored goal, for \bmGoals{} goals.
Deliberately \emph{not} set aside: the goals \beer{}'s encoder produced no script for, and the goals on which a solver was killed at the wall.
The first is encoder coverage---dominated by the files using Atelier~B's \texttt{Struct} records, a construct specific to Atelier~B rather than to the B language of the B-Book~\cite{abrial1996bbook}, which \beer{} does not aim to support---and the second is a failure to discharge within the budget; both are real outcomes of running the tool.

\paragraph*{Two scorings.}
Every result in this chapter uses one of two scoring rules, and confusing them changes every conclusion, so they are fixed here once.

\begin{description}[font=\normalfont\itshape]
  \item[strict] All \bmGoals{} goals; a goal counts as \emph{proved} iff cvc5 answered \smtinl{unsat}.
    A goal \beer{} could not encode, and a goal ppTrans burned the \(300\)\,s wall on, both count as failures to prove.
    This is what a user of either tool experiences end to end, and it is the headline scoring.
  \item[verdict-only] The \bmHead{} goals for which \emph{both} pipelines obtained a solver verdict (\smtinl{unsat} or \smtinl{unknown}).
    This isolates the quality of an encoding from the robustness and coverage of the encoder that produced it, at the cost of discarding \bmDropped{} goals---\(37\,\%\) of the corpus---and discarding them asymmetrically, since each side loses a different population.
\end{description}

\Cref{sec:beer-results-global} reports both scorings, because the gap between them is itself a result.
The structural analyses of \Cref{sec:beer-results-where,sec:beer-results-alone} are verdict-only throughout: they ask which encoding suits which kind of goal, and a goal that never reached the solver cannot answer that.
Every figure names its scoring, and the reference tables closing the chapter carry both side by side.
